% =============================================================================
%  «ТЕОРИЯ СИСТЕМ» — ствол серии «Архитектура Рассуждения»
%  Блок 5. ЦЕЛОЕ СВЕРХ ЧАСТЕЙ: эмерджентность.  Разделы 5.1–5.5 (черновик прозы).
%  Стиль: «выводим и показываем»; репо-якоря — в приложении, не в сносках.
%  Самокомпилируемо: pdflatex blok-5-celoe.tex (см. память latex-build-windows).
% =============================================================================
\documentclass[12pt,a4paper]{article}
\usepackage{cmap}
\usepackage[T2A]{fontenc}
\usepackage[utf8]{inputenc}
\usepackage[russian]{babel}
\usepackage{paratype}
\usepackage{amsmath,amssymb,amsthm}
\usepackage{listings}
\usepackage[expansion=false]{microtype}
\usepackage[margin=28mm]{geometry}

\newcommand{\rezultat}{\medskip\noindent\emph{Что отсюда следует.}\ }

\title{\textbf{Теория Систем}\\[2mm]\large Блок 5. Целое сверх частей: эмерджентность}
\author{}
\date{}

\begin{document}
\maketitle
\thispagestyle{empty}

\noindent\small\textit{Черновик: разделы 5.1--5.5. Продолжает Блок~4 (время). Берём \emph{много}
систем как целое и спрашиваем, что у целого есть \emph{сверх} частей. Машинные якоря --- в
Приложении.}\normalsize

\bigskip

\noindent Одну систему во времени мы разобрали. Возьмём теперь \emph{много} систем сразу.
\emph{Эмерджентность} --- это то, что есть у целого сверх его частей; и её можно определить
\emph{точно}. Она живёт на ярусе \emph{Roles}, а само соединение частей --- своя операция на каждом
ярусе.

% =============================================================================
\section*{5.1\quad Свойство и наблюдаемое}
% =============================================================================

\emph{Свойство} системы --- это функция от системы к «да\,/\,нет»; \emph{наблюдаемое} --- к значению
вообще. Они \emph{первоклассны}: о них можно говорить и доказывать наравне с самими системами.

Здесь же --- тихий, но важный факт. Всякая система удовлетворяет \emph{собственной} конституции (это
функциональный факт, §2.1): существующее \emph{познаваемо} --- сама его суть доступна. Это реализованная
познаваемость, и это мост к эпистемологии (\emph{быть значит быть познаваемым}).

И ещё одно --- \emph{карта часть$\to$целое}: правило, по которому из соотнесённости частей собирается
соотнесённость целого.

\rezultat
Свойства первоклассны, и есть карта часть$\to$целое. Теперь --- решающий вопрос: \emph{всякое ли}
свойство целого лежит в \emph{образе} этой карты?

% =============================================================================
\section*{5.2\quad Эмерджентность как свойство вне образа карты}
% =============================================================================

\textbf{Эмерджентное свойство --- то, которого нет в образе карты часть$\to$целое.} Его нельзя получить,
собрав свойства частей: оно принадлежит целому и только целому. Это и есть точное «сверх частей».

Перед нами не туман, а структурный факт: свойство целого либо лежит в образе карты (тогда оно
\emph{сводимо} к частям), либо вне его (тогда оно \emph{эмерджентно}). Третьего не дано.

\rezultat
Эмерджентность определена \textbf{точно и проверяемо}: свойство целого \emph{вне образа} карты
часть$\to$целое. Не метафора «больше суммы», а названное место --- дополнение к образу.

% =============================================================================
\section*{5.3\quad Роды и мера}
% =============================================================================

Эмерджентность \emph{таксономизируется}. Четыре страты, каждая со своим примером: \emph{сводимое}
(совпадает с произведением, лежит в образе); \emph{супер-аддитивное} (целое соотносит \emph{больше} ---
коллективная эмерджентность); \emph{суб-аддитивное} (целое соотносит \emph{меньше} --- корреляция,
ограничение); \emph{несепарабельное} (запутанность).

И эмерджентность \emph{измеряется}. На рациональных амплитудах есть детерминантный свидетель ---
\emph{рациональное} число, отделяющее сводимое от эмерджентного; а оптимум этой меры --- role-limit.

\rezultat
Эмерджентность не только \emph{опознаётся}, но и \emph{измеряется} --- и сепаратор рационален. Роды
названы; континуум появляется лишь в \emph{экстремуме} меры. (Честно: исчерпываемость родов не
заявлена --- это потребовало бы разрешимости.)

% =============================================================================
\section*{5.4\quad Запутанность как инстанс}
% =============================================================================

\emph{Запутанность} --- частный и важнейший случай: композит, который \emph{не} есть категорное
произведение. Её алгебраическое ядро --- чётность Bell\,/\,GHZ --- служит свидетелем несепарабельности.

Мера здесь конкретна. Детерминант матрицы амплитуд $2{\times}2$ --- рациональная «конкурренция»; детектор
\emph{полон и точен}: детерминант нулевой \emph{тогда и только тогда}, когда состояние факторизуемо. И
сепаратор, и его значение рациональны; \emph{только экстремум} --- амплитуды $1/\sqrt{2}$, граница
Цирельсона $2\sqrt{2}$ --- есть role-limit.

{\sloppy
Сверх того, системы образуют \emph{симметрично-моноидальную} структуру (тензор-бифунктор, единица,
симметрия) с \emph{дагером} на амплитудах. Важная локализация: запутанных \emph{состояний} на
реляционном ярусе нет --- всякое состояние произведения есть пара; запутанность живёт в
ролях\,/\,амплитудах. И всё замыкается двойственностью: \emph{произведение разделимо $\Leftrightarrow$
сумма несёт кросс-эмерджентность}.\par}

\rezultat
Запутанность --- \textbf{инстанс} общей эмерджентности, измеримый \emph{рационально}; континуум входит
лишь экстремумом ($\sqrt{2}$, Цирельсон). Полная дагер-компактность --- честная стена $\sqrt{2}$\,/\,Гильберт,
названная, а не обойдённая.

% =============================================================================
\section*{5.5\quad Исчисление соединения ярусов}
% =============================================================================

Как соединяются \emph{ярусы} при сборке целого? Каждый --- своей операцией \emph{возрастающей
связанности}. \emph{Elements} соединяются свободно --- произведение, безусловно. \emph{Roles} ---
встречей, и здесь возможна эмерджентность: целое может соотносить больше суммы частей. \emph{Rules} ---
гейт-решёткой: соединимы лишь согласованные правила.

\rezultat
Соединение --- не одна операция, а \emph{три}, по ярусам, возрастающей связанности: свободно
(Elements) $<$ встреча (Roles) $<$ гейт (Rules). Это \textbf{операционный смысл} ранговой асимметрии
(§2.4): ранг правил здесь развёрнут в исчисление. (Честно: свидетели конкретны, полнота решётки не
заявлена.)

\medskip
\noindent\emph{Переход.} Соединение по горизонтали разобрано. Следующий блок идёт по \emph{вертикали}
--- к башне уровней.

% =============================================================================
\appendix
\section*{Приложение к Блоку 5 — машинные якоря}
% =============================================================================
\small\raggedright\sloppy
\noindent Всё перечисленное машинно проверено; единственная аксиома ветви --- \texttt{classic} (L3).

\begin{itemize}\itemsep2pt
\item \textbf{§5.1.} \texttt{Property}, \texttt{Observable}, \texttt{system\_has\_essence}
(познаваемость), \texttt{part\_reducible} ${:=}$ \texttt{separable}.\\
\emph{Файл:} \texttt{ERRProperty.v}.

\item \textbf{§5.2.} эмерджентность как свойство вне образа карты --- \texttt{ERREmergenceSystem.v}.

\item \textbf{§5.3.} роды \texttt{reducible}/\texttt{super\_additive}/\texttt{sub\_additive}/\texttt{non\_separable}
(\texttt{ERREmergenceTaxonomy.v}); мера \texttt{q\_det} (\texttt{ERREmergenceMeasure.v}).

\item \textbf{§5.4.} \texttt{separable}, \texttt{parity\_roles} (\texttt{ERREntanglement.v});
мера \texttt{q\_det}, \texttt{det\_zero\_factorizable\_pivot}, свидетель \texttt{ent}
(\texttt{ERREntanglementMeasure.v}); \texttt{err\_tensor}, \texttt{state\_of\_product\_factors},
\texttt{dag} (\texttt{ERRMonoidal.v}); дуальность \texttt{ERRCoproductEmergence.v}.

\item \textbf{§5.5.} \texttt{elements\_combine\_free}, \texttt{roles\_super\_additive},
\texttt{cmeet} (гейт-решётка) --- \texttt{ERRCombinationCalculus.v}.
\end{itemize}
\normalsize

% --- Дальше: Блок 6 (ВЕРТИКАЛЬ) по тезисному плану в Книги/Теория Систем/. ---

\end{document}
